% K31-30, JC-003 y P01 del inventario REV05.
% Residencia: después de la derivación del registro K y de la incidencia Witt;
% antes de utilizar la lectura ponderada para la recuperación del registro.
% Propietarios: output/IMPLEMENTACION_K_SEMILLAS_INCIDENCIA_20260918/
% WeightedIncidenceRecovery.lean (íntegro) y
% COMPOSICION_GENERATIVA_Y_RECUPERACION_K.tex:257--306.
% Codificador: PAQUETE_ARTICULO_I_INTEGRACION_ESTADO_CAMPOS_20260919/
% lean/biblioteca/CoxeterNeighbor.lean:322--343;
% lean/biblioteca/SectorIncidenceData.lean:27--33.
% Estatuto: RESULTADO_RECUPERADO, explicitación íntegra de las matrices que
% la exposición anterior remitía al código. La prueba Lean es la ya existente.

\section[Recuperación por incidencias hexádicas]{Recuperación íntegra del registro mediante incidencias hexádicas ponderadas}
\label{rh:sec:recuperacion}

El registro dodecafásico generado conserva doce componentes ordenados y
un origen de fase. La incidencia de Paley--Witt proporciona una lectura
adicional de ese registro: cada soporte hexádico reúne seis de sus
componentes mediante una suma ponderada por la propia incidencia. La
reconstrucción que se desarrolla aquí recibe los valores de esas sumas
desde el mismo registro. Su objeto es probar que las lecturas conservan
todas las componentes y exhibir el operador inverso con coeficientes
racionales explícitos.

\subsection{Codificador de incidencia y doce testigos explícitos}

En el sistema de coordenadas de incidencia ya fijado, sea
\[
A_W=\begin{pmatrix}
0&1&1&1&1&1\\1&0&1&2&2&1\\1&1&0&1&2&2\\
1&2&1&0&1&2\\1&2&2&1&0&1\\1&1&2&2&1&0
\end{pmatrix}\in M_6(\mathbb F_3).
\]
Para $w\in\mathbb F_3^6$, el codificador es
$\operatorname{Enc}(w)=(w,wA_W)\in\mathbb F_3^{12}$ y su soporte
es
\[
H(w)=\{j\in\{1,\ldots,12\}:\operatorname{Enc}(w)_j\ne0\}.
\]
La elección de una posición emplea índices de uno a doce; las
implementaciones con índices desde cero se transportan mediante
$j\mapsto j+1$.

Los siguientes doce testigos pertenecen al dominio del codificador:
\[
\begin{array}{c|c@{\qquad}c|c}
i&w^{(i)}&i&w^{(i)}\\\hline
1&112210&7&122020\\
2&122101&8&112002\\
3&121200&9&111001\\
4&111100&10&110122\\
5&121012&11&101221\\
6&111020&12&011111.
\end{array}
\]
Definamos $H_i=H(w^{(i)})$ y
$B_{ij}=\mathbf1_{\{j\in H_i\}}$. La evaluación del codificador da
\begingroup\small\setlength{\arraycolsep}{3pt}
\setcounter{MaxMatrixCols}{12}
\begin{equation}
B=\begin{pmatrix}
1&1&1&1&1&0&0&0&0&0&0&1\\
1&1&1&1&0&1&0&0&0&0&1&0\\
1&1&1&1&0&0&1&0&1&0&0&0\\
1&1&1&1&0&0&0&1&0&1&0&0\\
1&1&1&0&1&1&0&0&0&1&0&0\\
1&1&1&0&1&0&1&0&0&0&1&0\\
1&1&1&0&1&0&0&1&1&0&0&0\\
1&1&1&0&0&1&1&1&0&0&0&0\\
1&1&1&0&0&1&0&0&1&0&0&1\\
1&1&0&1&1&1&0&0&1&0&0&0\\
1&0&1&1&1&1&0&1&0&0&0&0\\
0&1&1&1&1&1&1&0&0&0&0&0
\end{pmatrix}.
\label{rh:eq:incidencia}
\end{equation}
\endgroup
Cada fila contiene seis entradas iguales a uno. En consecuencia, cada
$H_i$ es una hexada de la imagen del codificador, y la matriz $B$ se
construye a partir de esos soportes.

\subsection{Inversa racional y unicidad de recuperación}

Sea la matriz entera
\begingroup\small\setlength{\arraycolsep}{3pt}
\setcounter{MaxMatrixCols}{12}
\begin{equation}
M=\begin{pmatrix}
7&1&8&-7&7&-1&-4&8&-7&3&3&-15\\
7&7&-4&-1&1&-7&8&8&-7&3&-15&3\\
1&7&8&-7&7&-7&8&-4&-1&-15&3&3\\
1&1&2&5&-5&-1&-4&-4&-1&3&3&3\\
1&-5&-4&-1&1&5&2&-4&-1&3&3&3\\
-5&1&-4&-1&1&-1&-4&2&5&3&3&3\\
-5&-11&2&5&-5&11&-10&2&5&3&3&3\\
-5&-5&-10&11&-11&5&2&2&5&3&3&3\\
-11&-5&2&5&-5&5&2&-10&11&3&3&3\\
-11&-11&-4&17&1&11&-10&-10&11&3&3&3\\
-11&1&-10&11&-11&17&-4&-10&11&3&3&3\\
1&-11&-10&11&-11&11&-10&-4&17&3&3&3
\end{pmatrix}.
\label{rh:eq:inversa-entera}
\end{equation}
\endgroup

\begin{theorem}[Recuperación exacta desde doce lecturas hexádicas]
\label{rh:thm:doce}
Para $v\in\mathbb Q^{12}$, sean
\[
\operatorname{Read}(v)=Bv,
\qquad
\operatorname{Rec}(y)=\frac1{18}My.
\]
Se verifican
\[
\operatorname{Rec}(\operatorname{Read}(v))=v,
\qquad
\operatorname{Read}(\operatorname{Rec}(y))=y.
\]
Por tanto, todo $y\in\mathbb Q^{12}$ admite un único registro
$v\in\mathbb Q^{12}$ con $\operatorname{Read}(v)=y$.
\end{theorem}
\begin{proof}
Las multiplicaciones enteras de las matrices
\eqref{rh:eq:incidencia} y \eqref{rh:eq:inversa-entera} dan
$MB=BM=18I_{12}$. Cada entrada diagonal resulta de un producto
escalar igual a dieciocho y cada entrada fuera de la diagonal resulta
de un producto escalar nulo. Al dividir por dieciocho se obtienen
las dos identidades inversas.

En particular,
$\operatorname{Rec}(\operatorname{Read}(v))=(MB/18)v=v$ y
$\operatorname{Read}(\operatorname{Rec}(y))=(BM/18)y=y$.
La segunda identidad prueba la existencia, mediante
$v=\operatorname{Rec}(y)$. Si $Bu=Bv$, aplicar $M/18$ a ambos
lados da $u=v$, lo que prueba la unicidad.
\end{proof}

La comprobación de $B$ desde el codificador y las dos multiplicaciones
exactas están formalizadas en los teoremas
\texttt{incidence\_from\_encoder}, \texttt{inverse\_left} y
\texttt{inverse\_right}. Las identidades de composición corresponden
a \texttt{recover\_read} y \texttt{read\_recover}.
Su contenido matemático se encuentra íntegramente en las matrices y
en la prueba anterior; el certificado formal registra la comprobación
de esas mismas operaciones.

\begin{corollary}[Separación por todas las incidencias ponderadas]
Si $u,v\in\mathbb Q^{12}$ satisfacen
\[
\sum_{j\in H(w)}u_j=\sum_{j\in H(w)}v_j
\]
para todo $w\in\mathbb F_3^6$ cuyo soporte codificado tiene seis
posiciones, entonces $u=v$.
\end{corollary}
\begin{proof}
La hipótesis se aplica en particular a los doce testigos expuestos.
Por la definición de $B$, sus doce igualdades constituyen $Bu=Bv$.
El teorema~\ref{rh:thm:doce} permite recuperar ambas entradas y
concluye $u=v$.
\end{proof}

\begin{example}[Recuperación de una componente periférica]
Para el vector canónico $v=e_{12}$, las únicas lecturas no nulas son
la primera y la novena:
\[
y=Be_{12}=(1,0,0,0,0,0,0,0,1,0,0,0)^{\mathsf T}.
\]
La suma de las columnas primera y novena de $M$ es $18e_{12}$.
Por consiguiente, $My/18=e_{12}$. Las dos sumas hexádicas distinguen
esta posición en la lectura conjunta porque las restantes componentes
se cancelan exactamente en el recuperador.
\end{example}

\subsection{Recuperador simétrico sobre las ciento treinta y dos hexadas}

El sistema de Steiner $S(5,6,12)$ ya producido por la incidencia
contiene $132$ hexadas. Denotemos su conjunto por $\mathcal H$ y
definamos
\[
\mathcal I:\mathbb Q^{12}\longrightarrow\mathbb Q^{\mathcal H},
\qquad
(\mathcal I v)_H=\sum_{j\in H}v_j.
\]
Este lector tiene redundancia: conserva las doce componentes mediante
$132$ sumas, mientras que el teorema anterior proporciona una base de
doce de esas sumas.

\begin{theorem}[Inversa sobre la imagen del lector completo]
\label{rh:thm:steiner}
Se cumple
\[
\mathcal I^{\mathsf T}\mathcal I=36I_{12}+30J_{12},
\]
donde $J_{12}$ tiene todas sus entradas iguales a uno. Para
$y=\mathcal I v$, el registro se recupera mediante
\begin{equation}
396v_i=11\sum_{H\ni i}y_H-5\sum_{H\in\mathcal H}y_H,
\qquad i=1,\ldots,12.
\label{rh:eq:recuperador}
\end{equation}
\end{theorem}
\begin{proof}
La propiedad $S(5,6,12)$ implica que cada posición pertenece a
$\binom{11}{4}/\binom{5}{4}=66$ hexadas. Cada pareja de posiciones
distintas pertenece a $\binom{10}{3}/\binom{4}{3}=30$ hexadas.
Por tanto, el producto de la matriz de incidencia transpuesta por
ella misma tiene diagonal $66$ y entradas extradiagonales $30$, que
es $36I+30J$.

Sea $S=\sum_jv_j$. Los dos conteos anteriores dan
\[
\sum_{H\ni i}y_H=66v_i+30\sum_{j\ne i}v_j=36v_i+30S,
\qquad
\sum_Hy_H=66S.
\]
Multiplicar la primera igualdad por once y restar cinco veces la
segunda elimina $S$ y produce $396v_i$. Esto demuestra
\eqref{rh:eq:recuperador} para toda lectura perteneciente a la
imagen de $\mathcal I$.
\end{proof}

\begin{corollary}[Componentes uniforme y transversal]
Sobre $\mathbb R^{12}$, el lector completo tiene valor singular
$\sqrt{396}$ en la dirección uniforme y valor singular $6$ en su
complemento ortogonal. Su inversor sobre la imagen tiene norma
operatoria $1/6$.
\end{corollary}
\begin{proof}
La matriz $J$ actúa por doce sobre $\mathbb R\mathbf1$ y por cero
sobre $\mathbf1^\perp$. Los autovalores de $36I+30J$ son, por
tanto, $396$ y $36$. Los valores singulares son sus raíces
positivas; el mayor valor singular del inversor es el recíproco
del menor de ellos.
\end{proof}

\subsection{Composición con el registro generado}

Aplicando estas identidades al registro entero dodecafásico $K$, ya
producido con su orden y origen, las lecturas $BK$ o $\mathcal I K$
permiten recuperar exactamente sus doce componentes. La interpretación
posicional posterior toma ese mismo registro en base mil; los
operadores de Hadamard, las diferencias y los lectores hexádicos son
realizaciones recuperables de una procedencia común.

El sistema de doce coordenadas racionales admite cualquier entrada
$y\in\mathbb Q^{12}$ y la reconstruye como un registro racional
único. El subdominio correspondiente a los registros HMT generados
conserva, además, sus condiciones de integridad, anchura de las
tríadas, origen, selección y frontera. Para el sistema redundante de
$132$ lecturas, la inversa actúa sobre la imagen de $\mathcal I$.
Estas precisiones tipan el contenido recuperado sin modificar la
generación antecedente del registro.
